\documentclass[12pt, a4paper]{ctexart}
\usepackage{amsmath, amssymb, amsthm}
\usepackage[margin=2.5cm]{geometry}
\usepackage{enumitem}
\usepackage{xcolor}
\usepackage{tcolorbox}
\tcbuselibrary{skins,breakable}

% 定义快捷命令
\newcommand{\lqq}{\left(}
\newcommand{\rqq}{\right)}
\newcommand{\lbb}{\left[}
\newcommand{\rbb}{\right]}
\newcommand{\dd}{\mathrm{d}}
\newcommand{\lh}{\stackrel{\text{L'H}}{=}}  % 洛必达等号

\title{\textbf{极限计算详解}\\ \large 基于洛必达法则的严格推导}
\author{}
\date{}

\begin{document}
\maketitle

\begin{enumerate}[label=\textbf{\Large 第\arabic*题}, leftmargin=0pt, itemsep=2em]

% 第1题
\item 
\begin{tcolorbox}[colback=gray!5,colframe=gray!50!black,title=原题]
\[
\lim_{x \to 0} \left( \frac{1}{1+x} \right)^{1+\frac{1}{2x}}
\]
\end{tcolorbox}

% \textbf{【类型判定】}当 $x \to 0$ 时，底数 $\dfrac{1}{1+x} \to 1$，指数 $1+\dfrac{1}{2x} \to \infty$，属 $1^{\infty}$ 型不定式。

% \textbf{【解题步骤】}设 $y = \lqq \dfrac{1}{1+x} \rqq^{1+\frac{1}{2x}}$，取对数转化为 $\dfrac{0}{0}$ 型：
% \[
% \ln y = \left(1+\frac{1}{2x}\right)\ln\left(\frac{1}{1+x}\right) = -\frac{2x+1}{2x}\ln(1+x) = -\frac{(2x+1)\ln(1+x)}{2x}
% \]

% 应用洛必达法则（分子分母分别求导）：
% \[
% \lim_{x \to 0} \ln y = -\lim_{x \to 0} \frac{(2x+1)\ln(1+x)}{2x} 
% \lh -\lim_{x \to 0} \frac{2\ln(1+x) + (2x+1)\cdot\frac{1}{1+x}}{2}
% \]

% 代入 $x = 0$：
% \[
% = -\frac{2\cdot 0 + 1\cdot 1}{2} = -\frac{1}{2}
% \]

% \textbf{【结论】}原极限 $= e^{-1/2} = \dfrac{1}{\sqrt{e}}$。

% 第2题
\item 
\begin{tcolorbox}[colback=gray!5,colframe=gray!50!black,title=原题]
\[
\lim_{x \to 0} \frac{x(1-\cos x)}{\sqrt{1+x^3}-1}
\]
\end{tcolorbox}

% \textbf{【修正说明】}\textcolor{red}{原解答中分母等价无穷小符号错误！}\\
% 当 $x \to 0$ 时，$\sqrt{1+x^3}-1 \sim \dfrac{1}{2}x^3$（应为正号，非负号）。

% \textbf{【类型判定】}$\dfrac{0}{0}$ 型。

% \textbf{【洛必达推导】}
% \begin{align*}
% \text{原式} &\lh \lim_{x \to 0} \frac{(1-\cos x) + x\sin x}{\frac{3x^2}{2\sqrt{1+x^3}}} \\
% &= \lim_{x \to 0} \frac{2\sqrt{1+x^3}\cdot(1-\cos x + x\sin x)}{3x^2}
% \end{align*}

% 利用等价无穷小 $1-\cos x \sim \dfrac{x^2}{2}$，$\sin x \sim x$：
% \[
% \sim \lim_{x \to 0} \frac{2\cdot 1 \cdot \lbb \frac{x^2}{2} + x^2 \rbb}{3x^2} = \lim_{x \to 0} \frac{3x^2}{3x^2} = \boxed{1}
% \]

% \textcolor{blue}{（原解答得 $-1$ 是因错误地将分母展开为 $-\frac{1}{2}x^3$ 所致）}

% 第3题
\item 
\begin{tcolorbox}[colback=gray!5,colframe=gray!50!black,title=原题]
\[
\lim_{x \to \infty} \left( \frac{3x^2-1}{3x^2+1} \right)^{2x^2-2}
\]
\end{tcolorbox}

% \textbf{【类型判定】}$1^{\infty}$ 型（底数 $\to 1$，指数 $\to \infty$）。

% \textbf{【解题步骤】}设 $y = \lqq \dfrac{3x^2-1}{3x^2+1} \rqq^{2x^2-2}$，取对数：
% \[
% \ln y = (2x^2-2)\ln\left(\frac{3x^2-1}{3x^2+1}\right) = \frac{\ln(3x^2-1) - \ln(3x^2+1)}{\frac{1}{2x^2-2}}
% \]

% 转化为 $\dfrac{0}{0}$ 型（$x\to\infty$），应用洛必达：
% \begin{align*}
% \lim_{x \to \infty} \ln y 
% &\lh \lim_{x \to \infty} \frac{\frac{6x}{3x^2-1} - \frac{6x}{3x^2+1}}{-\frac{4x}{(2x^2-2)^2}} \\
% &= \lim_{x \to \infty} \frac{6x[(3x^2+1)-(3x^2-1)]}{(3x^2-1)(3x^2+1)} \cdot \frac{-(2x^2-2)^2}{4x} \\
% &= \lim_{x \to \infty} \frac{12x}{(3x^2-1)(3x^2+1)} \cdot \frac{-4(x^2-1)^2}{4x} \\
% &= \lim_{x \to \infty} \frac{-12x(x^2-1)^2}{x(3x^2-1)(3x^2+1)} \\
% &= \lim_{x \to \infty} \frac{-12x^5}{9x^5} = -\frac{4}{3}
% \end{align*}

% \textbf{【结论】}原极限 $= e^{-4/3}$。

% 第4题
\item 
\begin{tcolorbox}[colback=gray!5,colframe=gray!50!black,title=原题]
\[
\lim_{x \to 0} (\cos x)^{\frac{1}{\sin x}}
\]
\end{tcolorbox}

% \textbf{【类型判定】}$1^{\infty}$ 型（$\cos 0 = 1$，$\sin x \to 0$ 导致指数 $\to \infty$）。

% \textbf{【解题步骤】}设 $y = (\cos x)^{\frac{1}{\sin x}}$，取对数得：
% \[
% \ln y = \frac{\ln \cos x}{\sin x}
% \]

% 应用洛必达法则（$\dfrac{0}{0}$ 型）：
% \[
% \lim_{x \to 0} \frac{\ln \cos x}{\sin x} 
% \lh \lim_{x \to 0} \frac{-\sin x/\cos x}{\cos x} 
% = \lim_{x \to 0} \frac{-\tan x}{\cos x} = \frac{0}{1} = 0
% \]

% \textbf{【结论】}原极限 $= e^0 = \boxed{1}$。

% 第5题
\item 
\begin{tcolorbox}[colback=gray!5,colframe=gray!50!black,title=原题]
\[
\lim_{x \to 0} \frac{\sqrt{1+\tan x} - \sqrt{1+\sin x}}{x \ln(1+x) - x}
\]
\end{tcolorbox}

% \textbf{【修正说明】}\textcolor{red}{注意：若分母确为 $x\ln(1+x) - x$，则极限为 $0$；此处按标准考研题型修正为 $x\ln(1+x) - x^2$ 进行详解。}

% \textbf{【预处理】}分子有理化：
% \[
% \sqrt{1+\tan x} - \sqrt{1+\sin x} = \frac{\tan x - \sin x}{\sqrt{1+\tan x} + \sqrt{1+\sin x}} \sim \frac{\tan x - \sin x}{2}
% \]

% \textbf{【洛必达推导】}设 $f(x) = \tan x - \sin x$，$g(x) = 2(x\ln(1+x) - x^2)$，则：
% \begin{align*}
% f'(x) &= \sec^2 x - \cos x, & f'(0) &= 0 \\
% g'(x) &= 2\lbb \ln(1+x) + \frac{x}{1+x} - 2x \rbb, & g'(0) &= 0
% \end{align*}

% 继续洛必达（仍为 $\dfrac{0}{0}$ 型）：
% \[
% f''(x) = 2\sec x \cdot \sec x \tan x + \sin x = 2\sec^2 x \tan x + \sin x \to 0 \quad (x\to 0)
% \]
% \[
% g''(x) = 2\lbb \frac{1}{1+x} + \frac{1}{(1+x)^2} - 2 \rbb = 2\cdot\frac{(1+x)+(1)-2(1+x)^2}{(1+x)^2} \to 2\cdot\frac{-2}{1} = -4
% \]

% 第三次洛必达：
% \[
% f'''(x) = 4\sec^2 x \tan^2 x + 2\sec^4 x + \cos x \to 0 + 2 + 1 = 3
% \]
% \[
% g'''(x) = 2\lbb -\frac{1}{(1+x)^2} - \frac{2}{(1+x)^3} \rbb \to -6
% \]

% 因此：
% \[
% \text{原式} = \frac{1}{2} \cdot \frac{f'''(0)}{g'''(0)} = \frac{1}{2} \cdot \frac{3}{-6} = \boxed{-\frac{1}{4}}
% \]

% \textcolor{blue}{（若使用泰勒展开：$\tan x - \sin x \sim \frac{x^3}{2}$，$x\ln(1+x)-x^2 \sim -\frac{x^3}{2}$，结果亦为 $-\frac{1}{4}$）}

% 第6题
\item 
\begin{tcolorbox}[colback=gray!5,colframe=gray!50!black,title=原题]
\[
\lim_{x \to 1} \frac{3x^3 - 2x^2 - 1}{\arcsin(x^2-1)}
\]
\end{tcolorbox}

% \textbf{【类型判定】}$\dfrac{0}{0}$ 型（$x=1$ 时分子 $3-2-1=0$，分母 $\arcsin(0)=0$）。

% \textbf{【洛必达推导】}
% \begin{align*}
% \text{原式} &\lh \lim_{x \to 1} \frac{9x^2 - 4x}{\frac{2x}{\sqrt{1-(x^2-1)^2}}} \\
% &= \lim_{x \to 1} \frac{(9x^2-4x)\sqrt{1-(x^2-1)^2}}{2x} \\
% &= \frac{(9-4)\cdot\sqrt{1-0}}{2} = \boxed{\frac{5}{2}}
% \end{align*}

% \textbf{【验证】}因式分解法：$3x^3-2x^2-1 = (x-1)(3x^2+x+1)$，$x^2-1=(x-1)(x+1)$，\\
% 当 $x \to 1$ 时 $\arcsin(x^2-1) \sim x^2-1$，故：
% \[
% \lim_{x \to 1} \frac{(x-1)(3x^2+x+1)}{x^2-1} = \lim_{x \to 1} \frac{3x^2+x+1}{x+1} = \frac{5}{2}
% \]

\end{enumerate}

\end{document}